\documentclass[a4paper,11pt]{article}
\usepackage[utf8]{inputenc}
\usepackage[T1]{fontenc}
\usepackage[ngerman]{babel}
\usepackage{amsmath,amssymb}
\usepackage{xcolor}
\usepackage{colortbl}
\usepackage{array}
\usepackage{tikz}
\usepackage{enumitem}
\usepackage[margin=2cm]{geometry}
\definecolor{bgdark}{HTML}{1E1E1E}
\definecolor{fgtext}{HTML}{E6E6E6}
\definecolor{gridcol}{HTML}{4A4A4A}
\definecolor{accent}{HTML}{6CB6FF}
\pagecolor{bgdark}
\color{fgtext}
\arrayrulecolor{fgtext}
\setlength{\parindent}{0pt}
\newcommand{\karo}[1]{\par\noindent\begin{tikzpicture}\draw[step=5mm,gridcol,very thin](0,0) grid (\linewidth,#1);\end{tikzpicture}\par\vspace{2mm}}
\newcommand{\aufgabe}[2]{\vspace{3mm}\par\noindent{\large\color{accent}\textbf{Aufgabe #1}}\hfill{\color{gridcol}\normalsize (#2 Punkte)}\par\vspace{1mm}}

\begin{document}
\begin{center}
  {\LARGE\color{accent}\textbf{Optimierungsverfahren}}\\[3pt]
  {\large Probeklausur 04 --- \textbf{Lineare Optimierung} (Vollabdeckung)}\\[2pt]
  {\small Deckt den kompletten Block \emph{Lineare Optimierung} ab: grafische Lösung \& Slackvariablen,
  Eckpunkte / Simplex (geometrisch), Standardform, Simplex-Tableau, ganzzahlige Optimierung.}
\end{center}
\vspace{2mm}
\noindent\textbf{Name:}\ \underline{\hspace{6cm}}\hfill\textbf{Datum:}\ \underline{\hspace{4cm}}\\[4pt]
\noindent\textbf{Gesamtpunkte:}\ 31 \hfill \textbf{Bearbeitungszeit:}\ ca.\ 45 Minuten \hfill \textbf{Hilfsmittel:}\ Open-Book, Taschenrechner\\[3pt]
{\color{gridcol}\rule{\linewidth}{0.4pt}}
\par\vspace{1mm}
\noindent{\small Begründen Sie Ihre Antworten kurz. Mathematik in sauberer Notation.}
\par\vspace{2mm}
\noindent\renewcommand{\arraystretch}{1.2}
\begin{center}
\begin{tabular}{|c|c|c|c|c|c|c|}
\hline
Aufgabe & 1 & 2 & 3 & 4 & 5 & $\Sigma$\\\hline
Punkte & 6 & 6 & 6 & 8 & 5 & 31\\\hline
erreicht & & & & & & \\\hline
\end{tabular}
\end{center}

% ===== 1 =====
\clearpage
\aufgabe{1 --- Grafische Lösung \& Slackvariablen}{6}
Gegeben sei das lineare Optimierungsproblem
\[
\max_{\vec x}\ f(\vec x)=3x_1+2x_2 \quad\text{u.d.N.}\quad
\underbrace{x_1+x_2\le 4}_{g_1},\ \ \underbrace{x_1+3x_2\le 6}_{g_2},\ \ x_1,x_2\ge 0.
\]
\begin{center}
\begin{tikzpicture}[scale=0.9]
  \fill[white,opacity=0.06] (0,0)--(4,0)--(3,1)--(0,2)--cycle;
  \draw[gridcol,very thin,step=1] (0,0) grid (6,5);
  \draw[->,fgtext] (0,0)--(6.4,0) node[right]{$x_1$};
  \draw[->,fgtext] (0,0)--(0,5.4) node[above]{$x_2$};
  \foreach \i in {1,2,3,4,5,6}{\node[fgtext,below] at (\i,0){\small\i};}
  \foreach \j in {1,2,3,4,5}{\node[fgtext,left] at (0,\j){\small\j};}
  \draw[accent,very thick] (4,0)--(0,4) node[above right,pos=0.12]{\small$g_1{:}\,x_1{+}x_2{=}4$};
  \draw[accent,very thick] (6,0)--(0,2) node[above right,pos=0.55]{\small$g_2{:}\,x_1{+}3x_2{=}6$};
\end{tikzpicture}
\end{center}
\begin{enumerate}[label=(\alph*),nosep,leftmargin=1.8em]
  \item (4\,P) Bestimmen Sie die optimale Ecke $\vec x^{\,*}$ und den optimalen Wert $f(\vec x^{\,*})$.
  \item (2\,P) Welche Werte nehmen die Slackvariablen $s_1,s_2$ im Optimum an? Welche Nebenbedingung ist aktiv?
\end{enumerate}
\karo{7cm}

% ===== 2 =====
\clearpage
\aufgabe{2 --- Eckpunkte \& Simplex-Suchverlauf (geometrisch)}{6}
Gegeben sei das lineare Optimierungsproblem
\[
\max_{\vec x}\ f(\vec x)=x_1+2x_2 \quad\text{u.d.N.}\quad
x_1\le 4,\ \ x_2\le 3,\ \ x_1+x_2\le 6,\ \ x_1,x_2\ge 0.
\]
\begin{center}
\begin{tikzpicture}[scale=0.8]
  \fill[white,opacity=0.06] (0,0)--(4,0)--(4,2)--(3,3)--(0,3)--cycle;
  \draw[gridcol,very thin,step=1] (0,0) grid (7,5);
  \draw[->,fgtext] (0,0)--(7.4,0) node[right]{$x_1$};
  \draw[->,fgtext] (0,0)--(0,5.4) node[above]{$x_2$};
  \foreach \i in {1,2,3,4,5,6,7}{\node[fgtext,below] at (\i,0){\small\i};}
  \foreach \j in {1,2,3,4,5}{\node[fgtext,left] at (0,\j){\small\j};}
  \draw[accent,very thick] (4,0)--(4,5) node[above]{\small$x_1{=}4$};
  \draw[accent,very thick] (0,3)--(7,3) node[right]{\small$x_2{=}3$};
  \draw[accent,very thick] (6,0)--(0,6) node[above right,pos=0.2]{\small$x_1{+}x_2{=}6$};
\end{tikzpicture}
\end{center}
\begin{enumerate}[label=(\alph*),nosep,leftmargin=1.8em]
  \item (3\,P) Bestimmen Sie rechnerisch \emph{alle} Eckpunkte des zulässigen Bereichs (auch über
        Schnittpunkte der Geraden) und werten Sie $f$ in jeder Ecke aus.
  \item (3\,P) Beschreiben Sie den Suchverlauf des Simplex-Verfahrens mit Start in $(0,0)$
        (entlang der Kante größter Verbesserung) und geben Sie $\vec x^{\,*}$ und $f(\vec x^{\,*})$ an.
\end{enumerate}
\karo{6.5cm}

% ===== 3 =====
\clearpage
\aufgabe{3 --- Standardform der Linearen Optimierung}{6}
Überführen Sie das folgende Problem in die \textbf{Standardform} (nur Minimierung, nur
Gleichungs-Nebenbedingungen, alle Variablen $\ge 0$, rechte Seiten $\ge 0$):
\[
\max_{\vec x}\ 2x_1+3x_2 \quad\text{u.d.N.}\quad
x_1+x_2\le 5,\ \ x_1-x_2\ge 1,\ \ x_1\ \text{frei},\ \ x_2\ge 0.
\]
Geben Sie alle nötigen Umformungen explizit an (Slack-/Surplus-Variablen, Aufspaltung der
freien Variablen, $\max\!\to\!\min$).
\karo{9cm}

% ===== 4 =====
\clearpage
\aufgabe{4 --- Simplex (Tableau-Verfahren)}{8}
Lösen Sie das folgende Problem mit dem \textbf{Simplex-Tableau} (Schlupfvariablen $s_1,s_2$):
\[
\max_{\vec x}\ f=3x_1+2x_2 \quad\text{u.d.N.}\quad
x_1+x_2\le 4,\ \ x_1+3x_2\le 6,\ \ x_1,x_2\ge 0.
\]
Geben Sie für jede Iteration Pivotspalte, Quotiententest, Pivotzeile und das aktualisierte
Tableau an. Nennen Sie am Ende $\vec x^{\,*}$ und $f^{*}$. (Vergleich mit Aufgabe~1 möglich.)
\karo{11cm}

% ===== 5 =====
\clearpage
\aufgabe{5 --- Ganzzahlige Optimierung}{5}
Gegeben sei
\[
\max_{\vec x}\ f=x_1+x_2 \quad\text{u.d.N.}\quad
2x_1+x_2\le 8,\ \ x_1+2x_2\le 8,\ \ x_1,x_2\ge 0,\ \ x_1,x_2\in\mathbb Z.
\]
\begin{enumerate}[label=(\alph*),nosep,leftmargin=1.8em]
  \item (2\,P) Bestimmen Sie die \emph{kontinuierliche} optimale Lösung und $f$.
  \item (3\,P) Bestimmen Sie die \emph{ganzzahlige} optimale Lösung. Liefert einfaches Abrunden der
        kontinuierlichen Lösung bereits das ganzzahlige Optimum?
\end{enumerate}
\karo{7cm}

\vfill
\begin{center}{\color{gridcol}\small--- Ende Probeklausur 04 (Lineare Optimierung, 31 Punkte) ---}\end{center}
\end{document}
